\documentclass[a4paper,11pt]{article}
\usepackage[utf8]{inputenc}
\usepackage[T1]{fontenc}
\usepackage[french]{babel}
\usepackage{geometry}
\geometry{left=2cm,right=2cm,top=1.5cm,bottom=2cm}
\usepackage{amsmath,amssymb}
\usepackage{enumitem}
\usepackage{array}
\usepackage{booktabs}
\usepackage{xcolor}
\usepackage{tcolorbox}
\usepackage{fancyhdr}
\usepackage{tikz}
\usepackage{lastpage}

% Couleurs
\definecolor{bleupro}{HTML}{0969da}
\definecolor{orange}{HTML}{d97706}
\definecolor{vertpro}{HTML}{2f855a}
\definecolor{rouge}{HTML}{c53030}
\definecolor{gris}{HTML}{718096}

% En-tête
\pagestyle{fancy}
\fancyhf{}
\lhead{\small CCF Physique-Chimie — Terminale ICCER}
\rhead{\small Fluides \& Courant alternatif}
\cfoot{\thepage/\pageref{LastPage}}
\renewcommand{\headrulewidth}{0.5pt}

% Boîtes
\tcbuselibrary{skins,breakable}

\newtcolorbox{formulairebox}{
  colback=yellow!5,
  colframe=orange,
  title={\textbf{FORMULAIRE — Rappels de cours}},
  fonttitle=\large\bfseries,
  breakable
}

\newtcolorbox{contextebox}{
  colback=blue!3,
  colframe=bleupro,
  leftrule=4pt,
  rightrule=0pt,
  toprule=0pt,
  bottomrule=0pt,
  boxsep=4pt,
  left=8pt
}

\newtcolorbox{reponsebox}[1][4cm]{
  colback=white,
  colframe=gris!40,
  boxrule=0.5pt,
  arc=3pt,
  height=#1,
  valign=top,
  before skip=4pt,
  after skip=8pt
}

\newcommand{\pts}[1]{\hfill\textcolor{gris}{\small\textbf{#1}}}
\newcommand{\comp}[1]{\colorbox{blue!10}{\scriptsize\textbf{#1}}}

\setlength{\parindent}{0pt}
\setlength{\parskip}{6pt}

\begin{document}

% ===== PAGE DE GARDE =====
\begin{center}
{\LARGE\bfseries\color{bleupro} Contrôle en Cours de Formation}\\[8pt]
{\large Physique-Chimie — Terminale Bac Pro ICCER — Groupement 1}\\[12pt]
{\Large\bfseries Débit de fluide \& Puissance en régime sinusoïdal}\\[20pt]

\begin{tabular}{|p{6cm}|p{6cm}|}
\hline
\textbf{Durée :} 45 minutes & \textbf{Barème :} /10 points \\
\hline
\textbf{Documents :} Calculatrice autorisée & \textbf{Formulaire :} ci-dessous \\
\hline
\textbf{Nom :} \dotfill & \textbf{Prénom :} \dotfill \\
\hline
\textbf{Classe :} \dotfill & \textbf{Date :} \dotfill \\
\hline
\end{tabular}
\end{center}

\vspace{10pt}

% ===== FORMULAIRE =====
\begin{formulairebox}

\subsection*{Débit et transport de fluide}

\begin{tabular}{>{\bfseries}p{4.5cm} p{5.5cm} p{4cm}}
\toprule
\textbf{Grandeur} & \textbf{Formule} & \textbf{Unités} \\
\midrule
Débit volumique & $Q_v = S \times v$ & m³/s, m², m/s \\[4pt]
Débit massique & $Q_m = \rho \times Q_v$ & kg/s \\[4pt]
Équation de continuité & $S_1 \times v_1 = S_2 \times v_2$ & \\[4pt]
Section d'un tuyau & $S = \pi \times \left(\dfrac{d}{2}\right)^2$ & m² \\[6pt]
Conversions & 1 m³/h $= \dfrac{1}{3\,600}$ m³/s & \\[4pt]
\bottomrule
\end{tabular}

\smallskip
$\rho_{\text{eau}} = 1\,000$ kg/m³ \quad Vitesse recommandée dans les canalisations : entre 0,5 et 1,5 m/s

\subsection*{Puissance en régime sinusoïdal}

\begin{tabular}{>{\bfseries}p{4.5cm} p{5.5cm} p{4cm}}
\toprule
\textbf{Grandeur} & \textbf{Formule} & \textbf{Unités} \\
\midrule
Puissance active & $P = U \times I \times \cos\varphi$ & W \\[4pt]
Puissance apparente & $S = U \times I$ & VA \\[4pt]
Puissance réactive & $Q = U \times I \times \sin\varphi$ & VAR \\[4pt]
Facteur de puissance & $\cos\varphi = \dfrac{P}{S}$ & (sans unité) \\[6pt]
Énergie & $E = P \times t$ & J (ou kWh) \\[4pt]
\bottomrule
\end{tabular}

\smallskip
\textbf{Triangle des puissances :} $S^2 = P^2 + Q^2$. Un $\cos\varphi$ proche de 1 est souhaitable (moins de pertes réactives).

\end{formulairebox}

\newpage

% ================================================================
\section*{Exercice 1 — TP numérique : dimensionner un réseau de chauffage \pts{5 points}}
% ================================================================

\begin{contextebox}
Un installateur thermique doit dimensionner le réseau d'eau chaude d'un bâtiment. Il dispose d'une \textbf{simulation numérique} pour tester différentes configurations avant de commander les canalisations.

\medskip
\textbf{Accès à la simulation :} \texttt{https://maths-sciences-lp.github.io/simulations/debit-fluide.html}

\medskip
\textbf{Rappel professionnel :} la vitesse de l'eau doit rester entre \textbf{0,5 m/s} (sinon dépôts) et \textbf{1,5 m/s} (sinon bruit et érosion).
\end{contextebox}

\vspace{6pt}
\textbf{Q1} \comp{APP} — \textbf{Manipulation :} ouvrir la simulation. Régler le débit à $Q_v = 1{,}8$ m³/h et le diamètre principal à $d_1 = 40$ mm. Relever les valeurs affichées et compléter le tableau. \pts{1 pt}

\begin{center}
\renewcommand{\arraystretch}{1.8}
\begin{tabular}{|>{\centering\arraybackslash}m{4cm}|>{\centering\arraybackslash}m{4cm}|}
\hline
\textbf{Grandeur} & \textbf{Valeur lue} \\
\hline
Section $S_1$ & \dotfill \\
\hline
Vitesse $v_1$ & \dotfill \\
\hline
Norme respectée ? & \dotfill \\
\hline
\end{tabular}
\renewcommand{\arraystretch}{1}
\end{center}

\textbf{Q2} \comp{REA} — \textbf{Vérification par le calcul :} retrouver la valeur de $v_1$ à partir de la formule $Q_v = S \times v$. Détailler les étapes. \pts{1 pt}

\begin{reponsebox}[3.5cm]
\dotfill

\dotfill

\dotfill

\dotfill
\end{reponsebox}

\textbf{Q3} \comp{ANA} — \textbf{Manipulation :} régler le diamètre des branches à $d_2 = 25$ mm. Relever la vitesse $v_2$ dans chaque branche. La norme est-elle respectée dans les branches ? \pts{1 pt}

\begin{reponsebox}[2.5cm]
$v_2$ lue dans la simulation = \dotfill m/s

Norme respectée ? \dotfill
\end{reponsebox}

\textbf{Q4} \comp{ANA} — \textbf{Manipulation :} la vitesse $v_1$ dans la canalisation principale est trop faible. En utilisant la simulation, faire \textbf{varier le diamètre $d_1$} jusqu'à obtenir une vitesse $v_1$ comprise entre 0,5 et 1,5 m/s. Compléter le tableau ci-dessous avec 3 essais. \pts{1 pt}

\begin{center}
\renewcommand{\arraystretch}{1.8}
\begin{tabular}{|>{\centering\arraybackslash}m{2.5cm}|>{\centering\arraybackslash}m{2.5cm}|>{\centering\arraybackslash}m{2.5cm}|>{\centering\arraybackslash}m{2.5cm}|}
\hline
\textbf{Essai} & \textbf{Diamètre $d_1$} & \textbf{Vitesse $v_1$} & \textbf{Conforme ?} \\
\hline
1 & 40 mm & \ldots\ldots m/s & \ldots\ldots \\
\hline
2 & \ldots\ldots mm & \ldots\ldots m/s & \ldots\ldots \\
\hline
3 & \ldots\ldots mm & \ldots\ldots m/s & \ldots\ldots \\
\hline
\end{tabular}
\renewcommand{\arraystretch}{1}
\end{center}

\textbf{Q5} \comp{COM} — Quel diamètre $d_1$ choisissez-vous pour la canalisation principale ? Justifier votre choix en tenant compte de la norme de vitesse et du coût (un tuyau plus petit coûte moins cher). \pts{1 pt}

\begin{reponsebox}[3.5cm]
Diamètre choisi : $d_1$ = \dotfill mm

Justification : \dotfill

\dotfill

\dotfill
\end{reponsebox}

\newpage

% ================================================================
\section*{Exercice 2 — Dimensionnement électrique d'une PAC \pts{5 points}}
% ================================================================

\begin{contextebox}
Un technicien installe une pompe à chaleur (PAC) air/eau pour chauffer un immeuble. La PAC est alimentée en monophasé 230 V / 50 Hz. D'après la plaque signalétique :
\begin{itemize}[nosep]
  \item Puissance active : $P = 3\,200$ W
  \item Facteur de puissance : $\cos\varphi = 0{,}85$
\end{itemize}
La PAC fonctionne 10 heures par jour pendant 180 jours (saison de chauffe). Le prix du kWh est 0,20 €.
\end{contextebox}

% Schéma triangle des puissances
\begin{center}
\begin{tikzpicture}[scale=0.6, every node/.style={font=\small}]
  \draw[thick, bleupro] (0,0) -- (7,0) node[midway, below=4pt] {$P$ (W) — active};
  \draw[thick, rouge] (7,0) -- (7,3.5) node[midway, right=4pt] {$Q$ (VAR)};
  \draw[thick, vertpro] (0,0) -- (7,3.5) node[midway, above=4pt, sloped] {$S$ (VA) — apparente};
  % Angle
  \draw (1.5,0) arc[start angle=0, end angle=26.6, radius=1.5];
  \node at (2,0.5) {$\varphi$};
\end{tikzpicture}
\end{center}

\textbf{Q1} \comp{REA} — Calculer la puissance apparente $S$ de la PAC. \pts{1 pt}

\begin{reponsebox}[3cm]
\dotfill

\dotfill

\dotfill
\end{reponsebox}

\textbf{Q2} \comp{REA} — Calculer l'intensité du courant $I$ absorbée par la PAC. \pts{1 pt}

\begin{reponsebox}[3cm]
\dotfill

\dotfill

\dotfill
\end{reponsebox}

\textbf{Q3} \comp{ANA} — Le disjoncteur de la ligne est calibré à 20 A. Est-il correctement dimensionné ? Justifier. \pts{1 pt}

\begin{reponsebox}[2.5cm]
\dotfill

\dotfill
\end{reponsebox}

\textbf{Q4} \comp{REA} — Calculer l'énergie consommée par la PAC sur une saison de chauffe (en kWh). En déduire le coût annuel. \pts{1 pt}

\begin{reponsebox}[3.5cm]
\dotfill

\dotfill

\dotfill

\dotfill
\end{reponsebox}

\textbf{Q5} \comp{COM} — Le propriétaire demande pourquoi sa facture est élevée. Le technicien remarque que le $\cos\varphi$ est de 0,85 au lieu de 0,95 (valeur idéale). Expliquer en quoi un mauvais facteur de puissance augmente l'intensité appelée, et pourquoi c'est un problème. \pts{1 pt}

\begin{reponsebox}[4cm]
\dotfill

\dotfill

\dotfill

\dotfill

\dotfill
\end{reponsebox}

\vfill

\begin{center}
\begin{tabular}{|c|c|c|c|}
\hline
\textbf{Exercice} & \textbf{Thème} & \textbf{Questions} & \textbf{Points} \\
\hline
1 & Débit de fluide (réseau chauffage) & Q1 à Q5 & /5 \\
\hline
2 & Puissance en régime sinusoïdal (PAC) & Q1 à Q5 & /5 \\
\hline
\multicolumn{3}{|r|}{\textbf{Total}} & \textbf{/10} \\
\hline
\end{tabular}
\end{center}

\end{document}
